\section{模型假设}
\begin{enumerate}[label=\arabic*.]
  \item \textbf{指标的相对可比性。}附件 A 的质量信号按原始定义取得；分类器 logits 经 softmax 后只用于文档之间的相对比较，不把模型输出概率解释为人工质量真值。方向、长度校正、缩尾与缺失填补参数只由 A1 确定，在 A2/A3 上冻结。
  \item \textbf{领域汇总与来源诊断。}用领域内文档主分的中位数表示该领域质量，以文档 bootstrap 描述给定样本下的波动。六来源分歧是主分算完后的描述性诊断，不参与评分、桥接或配比接口；A1 的 P95 阈值代表筛查预算，而非真实冲突率。
  \item \textbf{代理模型的有限迁移性。}在 A4/A5 训练的配比模型可用于 A6/A7 的同规模留出检验；向 60M、1B 的排序或绝对损失迁移须分别检验；尺度修正只用 60M 标定，1B 留作独立检验，不先验假定跨尺度迁移成立。A12--A15 为估算/外推数据，不作真实训练验证。
  \item \textbf{评价目标固定。}13 个验证域交叉熵等权汇总为主目标，评价权重不随优化的训练配比改变；pile\_cc 单域目标仅作副口径。候选配比的代理预测收益不等于重训后的实际收益。
  \item \textbf{跨附件质量锚点。}A 侧主质量分与 B 侧质量变量的映射参数不能从现有交叉熵水平直接识别；取同一数值口径作为主分析锚点属于建模选择，依赖质量聚合及 B 侧质量惩罚形式等条件，须做敏感性分析。
  \item \textbf{成本与时间情景。}训练、提质和注意力成本共用训练数据量 $D$；非规模技术进步以分析窗口内的连续时间项近似。该假设只支持条件预测，不涵盖新架构等突变。
\end{enumerate}
